% Part: lambda-calculus
% Chapter: introduction
% Section: primitive-recursion

\documentclass[../../../include/open-logic-section]{subfiles}

\begin{document}

\olfileid{lam}{int}{pr}
\olsection{توابع \usetoken{S}{lambda definable} تحت بازگشت اولیه بسته‌اند}

در مورد بازگشت اولیه سرانجام باید قدری کار کنیم و مرحله‌به‌مرحله پیش
برویم. مانند قبل، فرض کنید از پیش ترم‌های~$G'$ و~$H'$ را در اختیار داریم،
به‌گونه‌ای که هر یک به‌ترتیب تابع~$g$ و~$h$ را !!{lambda define}؛ می‌خواهیم
ترم~$F$ای بیابیم که تابع~$f$ با تعریف زیر را !!{lambda define}s:
\begin{align*}
f(0, \vec z) & = g(\vec z) \\
f(x+1, \vec z) & = h(x, f(x,\vec z), \vec z).
\end{align*}
پس به‌طور کلی، با داشتن ترم‌های لامبدای~$G'$ و~$H'$ کافی است ترم~$F$ای
بیابیم که
\begin{align*}
F(\num{0}, \vec z) & \equiv G(\vec z) \\
F(\overline{n+1}, \vec z) & \equiv H(\num{n}, F(\num{n}, \vec z), \vec z)
\end{align*}
را برای هر عدد طبیعی~$n$ برآورده کند؛ اینکه~$G'$ و~$H'$ هر یک به‌ترتیب
$g$ و~$h$ را !!{lambda define} بدین معناست که هرگاه عددنماهای~$\num{\vec m}$
را به‌جای~$\vec z$ بگذاریم، $F(\num{n+1}, \num{\vec m})$ به پاسخ درست
نرمال می‌شود.

اما برای این کار کافی است ترم~$F$ای بیابیم که
\begin{align*}
F(\num 0) & \equiv G \\
F(\overline {n+1}) & \equiv H(\num{n},F(\num{n}))
\intertext{را برای هر عدد طبیعی~$n$ برآورده کند، که در آن}
G & = \lambd[\vec z][G'(\vec z)] \text{ و}\\
H(u,v) & = \lambd[\vec z][H'(u,v(\vec z),\vec z)].
\end{align*}
به بیان دیگر، با ترفندهای لامبدایی می‌توانیم از نگرانی دربارهٔ پارامترهای
اضافی~$\vec z$ بپرهیزیم؛ این پارامترها در نمادگذاری لامبدا جذب می‌شوند.

پیش از تعریف ترم~$F$، به سازوکاری برای کار با زوج‌های مرتب نیاز داریم.
لم بعدی چنین سازوکاری را فراهم می‌کند.

\begin{lem}
ترم لامبدای~$D$ای وجود دارد که برای هر زوج از ترم‌های لامبدای~$M$ و~$N$،
$D(M,N)(\num{0}) \red M$ و~$D(M,N)(\num{1}) \red N$.
\end{lem}

\begin{proof}
نخست ترم لامبدای~$K$ را چنین تعریف کنید:
\[
K(y) = \lambd[x][y].
\]
به بیان دیگر، $K$ همان ترم~$\lambd[y][\lambd[x][y]]$ است. از دیدگاهی
دیگر، برای هر~$M$، عبارت~$K(M)$ تابع ثابتی است که با هر ورودی~$M$ را
برمی‌گرداند.

اکنون~$D(x,y,z)$ را با~$D(x,y,z) = z (K(y))x$ تعریف کنید. آنگاه داریم
\begin{align*}
D(M,N,\num 0) & \red \num 0 (K(N)) M \red M \text{ و}\\
D(M,N,\num 1) & \red \num 1 (K(N)) M \red K(N) M \red N,
\end{align*}
چنان‌که لازم بود.
\end{proof}

اندیشه این است که~$D(M,N)$ زوج~$\tuple{M, N}$ را بازنمایی می‌کند و اگر
فرض کنیم~$P$ چنین زوجی را بازنمایی می‌کند، $P(\num 0)$ و~$P(\num 1)$
به‌ترتیب افکنش‌های چپ و راست، $(P)_0$ و~$(P)_1$، را بازنمایی می‌کنند.
از این پس از همین نمادهای اخیر استفاده می‌کنیم.

\begin{lem}
توابع !!{lambda definable} تحت بازگشت اولیه بسته‌اند.
\end{lem}

\begin{proof}
باید نشان دهیم که با داشتن هر دو ترم~$G$ و~$H$، می‌توان ترم~$F$ای یافت
که
\begin{align*}
F(\num{0}) & \equiv G \\
F(\num{n+1}) & \equiv H(\num{n}, F(\num{n}))
\end{align*}
را برای هر عدد طبیعی~$n$ برآورده کند. اندیشهٔ کلی آن است که با استفاده
از عددنماها به‌عنوان تکرارگر، دنباله‌ای از \emph{زوج‌ها}
\[
\tuple{\num 0, F(\num 0)}, \tuple{\num 1, F(\num 1)}, \dots,
\]
را محاسبه کنیم. توجه کنید که زوج نخست همان~$\tuple{\num 0, G}$ است.
با داشتن زوج~$\tuple{\num{n}, F(\num{n})}$، زوج بعدی یعنی
$\tuple{\num{n+1}, F(\num{n+1})}$ باید با
$\tuple{\num{n+1}, H(\num{n}, F(\num{n}))}$ هم‌ارز باشد. ترم لامبدای~$T$
را چنان طراحی می‌کنیم که این گذار یک‌مرحله‌ای را انجام دهد.

جزئیات چنین است. $T(u)$ را با
\[
T(u) = \tuple{\fn{Succ}((u)_0), H((u)_0,(u)_1)}.
\]
تعریف کنید. اکنون به‌آسانی می‌توان بررسی کرد که برای هر عدد~$n$،
\[
T(\tuple{\num{n}, M}) \red \tuple{\num{n+1}, H(\num{n}, M)}.
\]
همان‌گونه که گفته شد، با داشتن~$G$ و~$H$، $F(u)$ را با
\[
F(u) = (u(T,\tuple{\num 0, G}))_1.
\]
تعریف کنید. به بیان دیگر، با ورودی~$\num{n}$، تابع~$F$، ترم~$T$ را
$n$ بار بر~$\tuple{\num 0, G}$ تکرار می‌کند و سپس مؤلفهٔ دوم را
برمی‌گرداند. برای آغاز داریم
\begin{enumerate}
\item $\num{0} (T, \tuple{\num 0, G}) \equiv \tuple{\num{0}, G}$
\item $F(\num{0}) \equiv G$
\end{enumerate}
با استقرا بر~$n$ می‌توان نشان داد که برای هر عدد طبیعی، موارد زیر برقرارند:
\begin{enumerate}
\item $\num{n+1}(T, \tuple{\num{0}, G}) \equiv \tuple{\num{n+1}, F(\num{n+1})}$
\item $F(\num{n+1}) \equiv H(\num{n}, F(\num{n}))$
\end{enumerate}
برای بند دوم داریم
\begin{align*}
F(\num{n+1}) & \red (\num{n+1}(T, \tuple{\num 0, G}))_1 \\
& \equiv (T(\num{n} (T, \tuple{\num 0, G})))_1 \\
& \equiv (T(\tuple{\num{n}, F(\num{n})}))_1 \\
& \equiv (\tuple{\num{n+1}, H(\num{n}, F(\num{n}))})_1 \\
& \equiv H(\num{n}, F(\num{n})).
\end{align*}
در سطر ماقبل آخر از فرض استقرا استفاده کرده‌ایم. برای بند نخست داریم
\begin{align*}
\num{n+1} (T, \tuple{\num 0, G}) &
\equiv T(\num{n} (T, \tuple{\num 0, G})) \\
& \equiv T( \tuple{\num{n}, F(\num{n})}) \\
& \equiv \tuple{\num{n+1}, H(\num{n}, F(\num{n}))} \\
& \equiv \tuple{\num{n+1}, F(\num{n+1})}.
\end{align*}
در سطر آخر از بند دوم استفاده کرده‌ایم. پس نشان دادیم که
$F(\num 0) \equiv G$ و برای هر~$n$، $F(\num {n+1}) \equiv H(\num
n, F(\num{n}))$؛ و این دقیقاً همان چیزی بود که نیاز داشتیم.
\end{proof}

\end{document}
